O texto indica que a força de arrasto $F_D$ depende de $V$ e também de três outras variáveis independentes: a largura da lâmina $d$, a densidade do fluido $\rho$ e a viscosidade do fluido $\mu$. Assim, uma representação bidimensional $F_D$ versus $V$ só faz sentido para condições em que $d$, $\rho$ e $\mu$ estão fixos.

Se escolhermos três valores diferentes para $\rho$ e três valores diferentes para $\mu$, teremos $3\times3=9$ combinações distintas de fluido. Para cada uma dessas 9 combinações, podemos ainda traçar 3 curvas diferentes correspondentes a 3 larguras de lâmina $d_1$, $d_2$ e $d_3$. Portanto:

\begin{itemize}
    \item 9 gráficos distintos correspondem às 9 combinações de ($\rho$, $\mu$);
    \item em cada gráfico há 3 curvas correspondentes às 3 larguras diferentes $d$.
\end{itemize}

Essa é a justificativa exata para a afirmação de que são necessários "nove (9) gráficos diferentes, cada um contendo três (3) curvas". Essa estrutura reflete a dependência tridimensional do problema antes da redução por análise dimensional.